\documentclass[12pt]{article}
\usepackage[utf8]{inputenc}
\usepackage[T1]{fontenc}
\usepackage[margin=1in]{geometry}
\usepackage{setspace}
\setlength{\parindent}{2em}
\setlength{\parskip}{1.25em}
\renewcommand{\baselinestretch}{1.0}
\makeatletter
\newcommand{\@subtitle}{}
\newcommand{\subtitle}[1]{\renewcommand{\@subtitle}{#1}}
\renewcommand{\@maketitle}{%
  \newpage
  \null
  \vskip 1em%
  \begin{center}%
    {\LARGE \@title \par}%
    \ifx\@subtitle\@empty\else
      \vskip 0.5em%
      {\large \@subtitle \par}%
    \fi
  \end{center}%
  \par
  \vskip 1em}
\makeatother
\usepackage{xcolor}
\usepackage[colorlinks=true,urlcolor=blue!55!black,linkcolor=black]{hyperref}

\title{\textit{Brown v. Board of Education (Brown I)}, 347 U.S. 483 (1954)}
\subtitle{Case Brief}
\author{}
\date{}

\begin{document}
\maketitle

\section{Memory Jogger}

In a unanimous landmark opinion authored by Chief Justice Earl Warren, the Supreme Court repudiated the ``separate but equal'' doctrine in public education, holding that state-sanctioned racial segregation of public schools is inherently unequal and violates the Fourteenth Amendment's Equal Protection Clause.

\section{Detailed Case Facts}

\textit{Brown v. Board of Education} consolidated four class-action lawsuits brought by Black schoolchildren and their legal representatives from Kansas, South Carolina, Virginia, and Delaware challenging state laws that required or permitted racial segregation in public schools. As detailed in Footnote 1 of the \href{https://www.courtlistener.com/opinion/105221/brown-v-board-of-education/}{CourtListener opinion}, \href{https://supreme.justia.com/cases/federal/us/347/483/}{Justia}, and \href{https://www.law.cornell.edu/supremecourt/text/347/483/USSC_PRO_347_483_1}{Cornell LII}, each case arose from distinct local conditions but presented the same core Fourteenth Amendment question:

\begin{itemize}
  \item \textbf{Kansas (\textit{Brown v. Board of Education of Topeka}):} Under Kan. Gen. Stat. \S{} 72-1724 (1949), cities with more than 15,000 residents were permitted, though not required, to maintain separate elementary schools for Black and white students. As documented by the \href{https://www.archives.gov/education/lessons/brown-v-board}{National Archives}, 13 Topeka parents (including Oliver Brown on behalf of his daughter Linda) sued on behalf of 20 children after nearby white elementary schools refused them admission. Because the federal district court found Topeka's Black and white schools substantially equal in buildings, transportation, curricula, and teacher qualifications, \textit{Brown} squarely challenged the validity of segregation itself.
  \item \textbf{South Carolina (\textit{Briggs v. Elliott}):} Black elementary and high school students in Clarendon County challenged S.C. Const. Art. XI, \S{} 7 and S.C. Code \S{} 5377 (1942), which mandated segregated public schools. Plaintiffs introduced groundbreaking psychological testimony from Dr. Kenneth Clark regarding the emotional harm of segregation.
  \item \textbf{Virginia (\textit{Davis v. County School Board of Prince Edward County}):} Black high school students challenged Va. Const. \S{} 140 and Va. Code \S{} 22-221 (1950), which required racially segregated schools, demonstrating stark disparities in physical plants, curricula, and transportation.
  \item \textbf{Delaware (\textit{Gebhart v. Belton}):} Black elementary and high school students in New Castle County challenged Del. Const. Art. X, \S{} 2 and Del. Rev. Code \S{} 2631 (1935), which mandated school segregation.
\end{itemize}

A fifth NAACP challenge from Washington, D.C., \href{https://supreme.justia.com/cases/federal/us/347/497/}{\textit{Bolling v. Sharpe}}, 347 U.S. 497 (1954), was argued alongside \textit{Brown} and decided the same day under the Fifth Amendment's Due Process Clause because the Fourteenth Amendment applies only to states.

Historically and politically, \textit{Brown} was the culmination of a decades-long litigation campaign by the NAACP Legal Defense and Educational Fund, led by Thurgood Marshall, Robert L. Carter, Spottswood W. Robinson III, Jack Greenberg, and Louis L. Redding. Since \textit{Plessy v. Ferguson}, 163 U.S. 537 (1896), ``separate but equal'' had legitimized Jim Crow segregation across the South. Beginning with graduate-school cases---\textit{Missouri ex rel. Gaines v. Canada} (1938), \textit{Sipuel v. Oklahoma} (1948), \textit{Sweatt v. Painter} (1950), and \textit{McLaurin v. Oklahoma State Regents} (1950)---the Court had chipped away at segregation by emphasizing both tangible and intangible inequalities without overruling \textit{Plessy}. At the same time, Cold War geopolitics heightened national scrutiny of American racial caste; the U.S. Department of Justice participated as \textit{amicus curiae} in both Terms urging invalidation of school segregation. After the cases were initially argued in December 1952 before a divided Court under Chief Justice Fred M. Vinson, the Court ordered reargument on five specific questions, including the original history of the Fourteenth Amendment. Vinson died in September 1953, and President Dwight D. Eisenhower appointed former California Governor Earl Warren as Chief Justice, who forged a unanimous Court before the decision was handed down on May 17, 1954.

\section{Procedural History}

\begin{itemize}
  \item \textbf{\textit{Brown} (Kansas):} A three-judge U.S. District Court (98 F. Supp. 797), convened under 28 U.S.C. \S{}\S{} 2281 and 2284, explicitly found that segregation had a detrimental effect on Black children, but denied relief under \textit{Plessy} because tangible facilities were substantially equal. Plaintiffs took a direct appeal under 28 U.S.C. \S{} 1253.
  \item \textbf{\textit{Briggs} (South Carolina):} A three-judge U.S. District Court (98 F. Supp. 529) upheld segregation over Judge Waties Waring's dissent but ordered equalization of inferior Black schools. After the Supreme Court vacated and remanded (342 U.S. 350), the district court found equalization largely achieved or underway and again denied admission to white schools (103 F. Supp. 920). Direct appeal followed under 28 U.S.C. \S{} 1253.
  \item \textbf{\textit{Davis} (Virginia):} A three-judge U.S. District Court (103 F. Supp. 337) ordered equalization of the inferior Black high school but upheld Virginia's segregation laws and denied immediate admission. Plaintiffs appealed directly under 28 U.S.C. \S{} 1253.
  \item \textbf{\textit{Gebhart} (Delaware):} The Delaware Court of Chancery (87 A.2d 862) ordered immediate admission of the Black plaintiffs to white schools due to tangible inequality (while also noting segregation's inherent harm). The Delaware Supreme Court affirmed (91 A.2d 137), intimating that defendants might seek modification once schools were equalized. State officials petitioned for certiorari, which the U.S. Supreme Court granted (344 U.S. 891).
  \item \textbf{U.S. Supreme Court:} Argued December 9--11, 1952; reargued December 7--9, 1953 (345 U.S. 972); decided May 17, 1954.
  \item \textbf{Disposition:} Held that racial segregation in public education deprives plaintiffs of the equal protection of the laws guaranteed by the Fourteenth Amendment; \textbf{restored the cases to the docket} for further argument on Questions 4 and 5 concerning the formulation of remedial decrees (later resolved in \href{https://supreme.justia.com/cases/federal/us/349/294/}{\textit{Brown II}}, 349 U.S. 294 (1955)).
\end{itemize}

\section{Judicial Votes}

\begin{itemize}
  \item \textbf{Majority (9 Justices -- Unanimous):} Chief Justice Earl Warren (author), joined by Justices Hugo L. Black, Stanley F. Reed, Felix Frankfurter, William O. Douglas, Robert H. Jackson, Harold H. Burton, Tom C. Clark, and Sherman Minton.
  \item \textbf{Concurrences:} None.
  \item \textbf{Dissents:} None.
\end{itemize}

\section{Holding}

\begin{itemize}
  \item \textbf{Judgment:} State laws requiring or permitting racial segregation in public schools violate the Equal Protection Clause of the Fourteenth Amendment; the cases were restored to the docket for remedial reargument.
  \item \textbf{Rule Announced:} In the field of public education, the doctrine of ``separate but equal'' has no place. Separate educational facilities are inherently unequal, and segregating public school children solely on the basis of race deprives minority children of equal educational opportunities even when physical facilities and other tangible factors are equal.
  \item \textbf{Holding vs. Dicta:} The holding is confined to \textit{de jure} racial segregation in public education under the Fourteenth Amendment's Equal Protection Clause. The Court expressly declined to decide whether school segregation also violates the Fourteenth Amendment's Due Process Clause, and it did not formally overrule \textit{Plessy v. Ferguson} outside the educational context in \textit{Brown I} itself (though subsequent \textit{per curiam} rulings rapidly extended \textit{Brown} to public beaches, buses, golf courses, and parks). Chief Justice Warren's discussion of the modern civic importance of public education and the inconclusiveness of the 1868 historical record explains the Court's interpretive method rather than establishing an independent right to education, and Footnote 11's citation of psychological studies served as corroborating support for the Kansas district court's factual finding rather than a factual prerequisite for future equal-protection challenges.
\end{itemize}

\section{Analysis of Opinions}

\subsection{Unanimous Opinion of the Court (Chief Justice Warren)}

Chief Justice Warren deliberately wrote a concise, accessible, non-accusatory opinion designed to be read by ordinary citizens and to minimize regional backlash. He began by addressing the extensive historical reargument on the adoption of the Fourteenth Amendment in 1868. Surveying congressional debates, state ratifications, and mid-nineteenth-century schooling practices, Warren concluded that the historical sources were ``[a]t best, \ldots inconclusive.'' Public education in 1868 was rudimentary in the North, barely established in the South, and practically nonexistent for Black Americans; compulsory attendance laws were virtually unknown. Accordingly, Warren declared that the Court ``cannot turn the clock back to 1868 when the Amendment was adopted, or even to 1896 when \textit{Plessy v. Ferguson} was written,'' and instead had to evaluate public education ``in the light of its full development and its present place in American life throughout the Nation.''

Warren next emphasized the central role of public schooling in modern democratic life: ``Today, education is perhaps the most important function of state and local governments.'' It is the foundation of good citizenship, cultural awakening, professional preparation, and normal social adjustment, such that ``where the state has undertaken to provide it,'' education ``is a right which must be made available to all on equal terms.''

Because lower courts in the consolidated cases had found that tangible factors were equal or being equalized, Warren framed the decisive question as whether racial segregation \textit{itself} deprives minority children of equal educational opportunities. Drawing on the intangible-factors reasoning of \textit{Sweatt v. Painter} and \textit{McLaurin v. Oklahoma State Regents}, Warren reasoned that intangible harms apply ``with added force to children in grade and high schools.'' Separating children solely because of race ``generates a feeling of inferiority as to their status in the community that may affect their hearts and minds in a way unlikely ever to be undone.'' Adopting the Kansas district court's finding that state-sanctioned segregation denotes inferiority and retards educational development, Warren stated in Footnote 11 that this finding was ``amply supported by modern authority''---citing Kenneth Clark's doll studies, Gunnar Myrdal's \textit{An American Dilemma}, and other social-science works---and expressly rejected any contrary language in \textit{Plessy}. Concluding that ``[s]eparate educational facilities are inherently unequal,'' the Court held that segregation denied plaintiffs equal protection and restored the cases to the docket to address the complex remedial questions.

\subsection{Concurrences and Dissents}

There were no separate concurrences or dissents. As noted in \href{https://supreme.justia.com/cases/federal/us/347/483/}{Justia's case commentary}, unanimity was not inevitable: Justices Reed and Jackson had initially considered writing separately or dissenting, and Justice Frankfurter had urged reargument partly to buy time to build consensus. Chief Justice Warren and his colleagues recognized that a fractured Court would invite defiant resistance across the segregated South, and their unified 9--0 voice became one of the decision's defining institutional strengths.

\section{Examples: Future Applications}

\begin{enumerate}
  \item \textbf{Same-Side Application (State-Mandated Segregation in Other Public Facilities):} A city enacts an ordinance requiring separate public parks, beaches, or municipal buses for Black and white residents, arguing that the facilities are physically identical and that \textit{Brown} was limited to public schools. Applying \textit{Brown}'s core anti-subordination and anti-classification principles, courts invalidate the ordinance as an unconstitutional badge of racial inferiority---just as the Supreme Court did in a series of post-\textit{Brown} \textit{per curiam} orders.
  \item \textbf{Same-Side Application (Racially Segregated Public Higher Education):} A state maintains a state-funded flagship university exclusively for white students and a newly built, physically comparable university exclusively for Black students. Under \textit{Brown} (building on \textit{Sweatt}), tangible parity does not cure state-enforced racial separation; separate public educational institutions remain inherently unequal under the Equal Protection Clause.
  \item \textbf{Opposite-Side Application (\textit{De Facto} Residential Imbalance Without State Action):} Two public elementary schools in a northern city have heavily imbalanced racial demographics solely because of private residential housing choices and neutral, race-blind neighborhood attendance zones, with no history of intentional state-mandated segregation. Because \textit{Brown} addressed \textit{de jure} segregation imposed ``with the sanction of law,'' purely private \textit{de facto} demographic imbalance does not by itself establish a Fourteenth Amendment violation without proof of discriminatory state intent.
  \item \textbf{Opposite-Side Application (Private, Non-State Religious School):} A strictly private religious academy that receives no state authority and is not a state actor adopts admissions policies challenged solely under the Fourteenth Amendment's Equal Protection Clause. Because the Fourteenth Amendment constrains only state action, \textit{Brown} does not directly apply to purely private institutions (though federal civil rights statutes or tax-exemption rules, as in \textit{Bob Jones University v. United States}, may still restrict them).
  \item \textbf{Fence-Sitter Application (Voluntary Race-Conscious School Integration Plans):} A school district that was never subject to a desegregation decree voluntarily uses individual students' race as a tiebreaker in school assignments to prevent racial isolation and promote integration. This scenario divides jurists over \textit{Brown}'s true meaning: an ``anti-classification'' reading (adopted by the plurality in \textit{Parents Involved in Community Schools v. Seattle School District No. 1}, 551 U.S. 701 (2007)) holds that \textit{Brown} forbids assigning any student to a school based on race, whereas an ``anti-subordination'' reading (advanced by Justice Breyer's dissent) argues that \textit{Brown} targeted caste-enforcing exclusion and permits voluntary measures to achieve integrated schools.
\end{enumerate}

\section{Critique}

Although \textit{Brown v. Board of Education} is universally revered as a canonical triumph of American constitutional law, its legal reasoning has sparked some of the most searching debates in constitutional theory.

From a \textbf{legal-process perspective}, Columbia Law professor Herbert Wechsler famously argued in \textit{Toward Neutral Principles of Constitutional Law} (1959)---analyzed in \href{https://scholarship.law.columbia.edu/cgi/viewcontent.cgi?article=3132&context=faculty_scholarship}{Columbia Law scholarship} and the \href{https://direct.mit.edu/ajle/article/doi/10.1162/ajle_a_00068/125059/WHAT-BROWN-TAUGHT-US-ABOUT-THE-CONSTITUTION}{\textit{American Journal of Law and Equality}}---that although he abhorred segregation, the Court's opinion failed to articulate a neutral constitutional principle. Wechsler framed segregation as a clash between Black citizens' freedom to associate and white citizens' freedom \textit{not} to associate, and questioned why the Court's reliance on Footnote 11's contested social-science studies could justify preferring one associational claim over the other. Yale professor Charles Black delivered the classic rejoinder, contending that Wechsler ignored the unmistakable social reality of Jim Crow: segregation was not a neutral denial of mutual association, but a system explicitly designed to stamp Black Americans as an inferior caste.

From a \textbf{progressive and critical race perspective}, Derrick Bell argued in \href{https://yale.imodules.com/s/1667/images/gid6/editor_documents/bellinterest.pdf}{\textit{Brown v. Board of Education and the Interest-Convergence Dilemma}} (1980) that the Court achieved unanimity in \textit{Brown} largely because Black Americans' quest for racial equality temporarily converged with the Cold War foreign-policy and economic interests of white elites seeking to counter Soviet propaganda abroad. Progressive critics also faulted the Court's reliance on psychological-stigma findings---which risked framing Black schools as inherently defective rather than condemning white supremacy directly---and criticized the decision to defer remedies to \href{https://supreme.justia.com/cases/federal/us/349/294/}{\textit{Brown II}}'s vague ``with all deliberate speed'' formula, which enabled over a decade of ``Massive Resistance'' across the South.

From an \textbf{originalist and textualist perspective}, \textit{Brown} has long served as the ultimate litmus test of interpretive legitimacy. Early original-intent theorists such as Raoul Berger argued that because the Thirty-Ninth Congress funded segregated schools in the District of Columbia and most ratifying states maintained segregated schools in 1868, Warren's explicit refusal to ``turn the clock back to 1868'' demonstrated that \textit{Brown} departed from original intent in favor of living constitutionalism. Modern originalists, however, have vigorously defended \textit{Brown}'s result while critiquing Warren's non-originalist opinion. Most notably, Michael W. McConnell argued in \href{https://chicagounbound.uchicago.edu/cgi/viewcontent.cgi?article=12612&context=journal_articles}{\textit{The Originalist Case for Brown v. Board of Education}} and his \href{https://chicagounbound.uchicago.edu/cgi/viewcontent.cgi?article=12613&context=journal_articles}{reply to Michael Klarman} that large majorities of the Reconstruction Congresses debating the Civil Rights Act of 1875 concluded that the Fourteenth Amendment forbade \textit{de jure} school segregation. Similarly, Robert Bork and Steven Calabresi have defended \textit{Brown} on the ground that the Equal Protection Clause's original public meaning established a general anti-caste principle incompatible with racial separation once segregation's actual inequality was recognized---even as critics like Michael Klarman contend that such originalist reconstructions stretch the historical record.

\section{Key Quotations}

\begin{quote}
In approaching this problem, we cannot turn the clock back to 1868 when the Amendment was adopted, or even to 1896 when Plessy v. Ferguson was written. We must consider public education in the light of its full development and its present place in American life throughout the Nation.
\end{quote}
\noindent --- Chief Justice Earl Warren, writing for a unanimous Court, 347 U.S. at 492--93 (\href{https://www.law.cornell.edu/supremecourt/text/347/483/USSC_PRO_347_483_1}{Cornell LII}; \href{https://www.courtlistener.com/opinion/105221/brown-v-board-of-education/}{CourtListener}).

\begin{quote}
Today, education is perhaps the most important function of state and local governments. Compulsory school attendance laws and the great expenditures for education both demonstrate our recognition of the importance of education to our democratic society. It is required in the performance of our most basic public responsibilities, even service in the armed forces. It is the very foundation of good citizenship. Today it is a principal instrument in awakening the child to cultural values, in preparing him for later professional training, and in helping him to adjust normally to his environment. In these days, it is doubtful that any child may reasonably be expected to succeed in life if he is denied the opportunity of an education. Such an opportunity, where the state has undertaken to provide it, is a right which must be made available to all on equal terms.
\end{quote}
\noindent --- Chief Justice Earl Warren, writing for a unanimous Court, 347 U.S. at 493 (\href{https://supreme.justia.com/cases/federal/us/347/483/}{Justia}).

\begin{quote}
To separate them from others of similar age and qualifications solely because of their race generates a feeling of inferiority as to their status in the community that may affect their hearts and minds in a way unlikely ever to be undone.
\end{quote}
\noindent --- Chief Justice Earl Warren, writing for a unanimous Court, 347 U.S. at 494 (\href{https://supreme.justia.com/cases/federal/us/347/483/}{Justia}).

\begin{quote}
We conclude that in the field of public education the doctrine of ``separate but equal'' has no place. Separate educational facilities are inherently unequal. Therefore, we hold that the plaintiffs and others similarly situated for whom the actions have been brought are, by reason of the segregation complained of, deprived of the equal protection of the laws guaranteed by the Fourteenth Amendment.
\end{quote}
\noindent --- Chief Justice Earl Warren, writing for a unanimous Court, 347 U.S. at 495 (\href{https://www.law.cornell.edu/supremecourt/text/347/483/USSC_PRO_347_483_1}{Cornell LII}).

\end{document}
